\documentclass[a4paper, 14pt]{extarticle}

\usepackage[T2A]{fontenc}
\usepackage[utf8]{inputenc}
\usepackage[russian]{babel}



\begin{document}
\setlength{\emergencystretch}{1em}
\begin{titlepage}
  \centering

  Общероссийская общественно-государственная организация 
  ``Всероссийское общество изобретателей и рационализаторов'' (ВОИР)                        

  \vfill

  {\Large\bfseries ТЕХНИЧЕСКОЕ ЗАДАНИЕ\par}
  \vspace{1cm}
  {\large На беспилотную авиационную систему (БАС)\par}
  
  \vfill

  \begin{flushright}
    Разработал: AIRES
  \end{flushright}

  \vfill
  
  Новосибирск\\
  \the\year
\end{titlepage}

\tableofcontents

\newpage

\section{Введение}
Настоящее техническое задание определяет требования к разработке аппаратной платформы беспилотной авиационной системы (БАС) квадрокоптерного типа, включающей контрукцию, силовую установку, систему питания, полётный контроллер, одноплатный компьютер и дополнительное оборудование, и предназначенной для выполнения захвата и переноса грузов-моделей в условиях макета ``robotrio'' посредством компьютерного зрения. 

Цель разработки -- создать систему, обеспечивающую выполнение полётного задания в рамках транспортной миссии в пределах макета. Эксплуатация предполагается в помещении. 

Данный документ устанавливает требования к БАС.

\section{Основание для разработки}

\section{Требование к БАС}

\subsection{Требования к функциональным характеристикам}

Разрабатываемая БАС должна обеспечивать:

\begin{itemize}
  \item взлёт со стартовой площадки и посадку на неё по команде оператора,
  \item перемещение в пределах макета по командам ручного управления и через указание целевого положения в координатах системы и(или) целевого вектора скорости,
  \item удержание позиции с точностью до 1 см по трём осям,
  \item приём команд, передача телеметрии, изображений камер и показаний датчиков наземной станции управления.
\end{itemize}

\subsection{Требования к техническим характеристикам}

\begin{itemize}
  \item Габариты БАС не должны превышать $30 \textit{см}$ в диаметре в боевом положении.
  \item Грузоподъёмность БАС должна быть не менее $500 \textit{г}$ полезной нагрузки (без учёта полётного контроллера).
\end{itemize}

\subsection{Требования к составу и конструкции}

Основное оборудование:

\begin{itemize}
  \item полётный контроллер с поддержкой автопилота PX4-Autopilot либо аналогов,
  \item система связи с наземной станцией, способной передавать сообщения Protobuf,
  \item крепления компонент,
  \item защита проводов,
  \item защита на винты.
\end{itemize}

Разрабатываемая БАС должна иметь возможность установки следующего дополнительного оборудования:

\begin{itemize}
  \item одноплатный компьютер Orange Pi5 Plus;
  \item камера USB+HDMI x3;
  \item лидар UART;
  \item система захвата груза-модели (манипулятор или электромагнит)
\end{itemize}

\subsection{Требования к электропитанию}

\begin{itemize}
  \item Обеспечить беcперебойное питание полётного контроллера и дополнительного оборудования.
  \item Аккумулятор должен быть сменным.
  \item Продолжительность полёта на одном заряде аккумулятора при полной нагрузке должна составлять не менее 15 минут. 
\end{itemize}

\subsection{Требования к безопасности}

\begin{itemize}
  \item Требуется наличие защитных корпусов на винтах БАС.
  \item Требуется блокировка случайного запуска системы.
  \item В случае отключения управляющего сигнала система должна немедленно перейти в безопасный режим и совершить посадку.
\end{itemize}

\subsection{Условия эксплуатации}

Работа системы предполагается в следующих условиях:

\begin{itemize}
  \item Помещение или защитный купол.
  \item Размеры макета - 3 на 2 метра.
\end{itemize}

\end{document}